\pdfoutput=1
\documentclass[11pt,a4paper]{article}
\usepackage{abhijit-research}
\renewcommand{\PaperCode}{P11}
\renewcommand{\PaperKind}{RESEARCH PAPER}
\renewcommand{\PaperVersion}{VERSION 1.0}
\renewcommand{\PaperDate}{AUGUST 2026}
\renewcommand{\PaperTitle}{Reenactable Closure}
\renewcommand{\PaperSubtitle}{Conditional Quantum-Causal Representation}
\renewcommand{\PaperFullTitle}{Reenactable Closure and Conditional Quantum-Causal Representation}
\renewcommand{\PaperShortTitle}{Reenactable Closure}
\renewcommand{\PaperKeywords}{formation calculus; fixed points; generated probes; quantum reconstruction; causal geometry; block connection}
\renewcommand{\PaperClassification}{Reenactment, probe closure, block scaling, trace exchange, and isospectrality are formal theorems. Quantum and geometric results are conditional representations. No new physical conversion law is claimed without an independently fixed observable.}
\renewcommand{\PaperCitation}{Singh, Abhijit. 2026. \textit{Reenactable Closure and Conditional Quantum-Causal Representation}. Research manuscript, version 1.0.}
\begin{document}
\MakeResearchTitle
\MakeResearchContents
\MakeResearchAbstract{%
A formation-first calculus is developed around an occurrence in which an active role, stabilized content, and successor standpoint arise together. Stabilized seats can reenact formation patterns; diagonal reenactment yields fixed points for seat-expressible transformers. Internally generated probes define complete signatures, equivalence, laws, and an idempotent completion, replacing an external meta-level closure. The invariant algebra supports sequential and unresolved-alternative composition. Under explicit reconstruction assumptions, one quotient admits finite complex quantum representation and another admits causal/Lorentzian geometry. A block connection couples active, stabilized, and conversion sectors. Its curvature forces first-order off-diagonal conversion and second-order diagonal backreaction; the quadratic terms obey trace exchange, and reciprocal maps yield identical nonzero singular spectra in the two sectors. Unsupported mixed-path physical interpretations are retired while the exact algebra is retained.
}
\section{Act, record, and changed standpoint}
The foundational asymmetry is operational: an occurring act is not identical to the stabilized record produced by that act, and stabilization changes the standpoint from which later operations proceed. A conclusion can be transmitted as content; the operation that made it meaningful must be reenacted.

Represent one primitive formation by
\begin{equation}
\Phi=\langle S\mid a\triangleright c\mid S^+\rangle,
\end{equation}
where $a$ is the active role, $c$ the stabilized role, $S$ the operative standpoint, and $S^+$ the successor standpoint. These are roles within one occurrence rather than independently fundamental objects.

\section{Seats and reenactment}
A seat $[P]$ is stabilized content capable of regenerating a formation pattern $P$. Reenactment $\uparrow[P]$ is successful when it reproduces the structural-functional invariant of $P$.

Let $T$ be any transformer expressible through seats. Define diagonal reenactment
\begin{equation}
\Delta(x)=\uparrow(x\bullet x).
\end{equation}
Construct a seat $q_T$ such that reenacting it on $x$ performs $T(\Delta(x))$, and put $\Omega_T=\Delta(q_T)$.

\begin{theorem}[Reenactment fixed point]
For every seat-expressible transformer $T$,
\begin{equation}
\Omega_T\simeq T(\Omega_T).
\end{equation}
\end{theorem}
\begin{proof}
By construction, reenacting $q_T$ on $x$ gives $T(\Delta(x))$. Substitution $x=q_T$ gives the result.
\end{proof}
If completion itself is seat-expressible, there exists a pattern invariant under completion. Finality means closure of representational form relative to generated probes, not cessation of events or discoveries.

\section{Generated probes and complete signatures}
For a pattern $P$, let $\mathcal P(P)$ be the least probe family closed under enactment, stabilization, successor-state change, composition, and available comparisons. Define the complete signature
\begin{equation}
\Sigma(P)=\bigl(p(P)\bigr)_{p\in\mathcal P(P)}.
\end{equation}
Patterns are equivalent when their complete generated signatures agree. A law is an automorphism preserving the signature and operation structure.

\begin{theorem}[Endogenous completion]
If the generated probe algebra is closed under composition, adjoining all generated probe responses is idempotent:
\begin{equation}
K(K(P))\simeq K(P).
\end{equation}
Law generation is likewise idempotent relative to the closed probe algebra.
\end{theorem}
This closes a recurring meta-level gap: probes are not imported by an unexamined external observer; they are generated from the enactment structure being completed.

\section{Formation invariant algebra}
Let $I(P)$ be the reenactment-equivalence class of $P$. Sequential formation and unresolved-alternative formation induce operations
\begin{equation}
I(P\odot Q)=I(P)\odot I(Q),\qquad I(P\boxplus Q)=I(P)\boxplus I(Q).
\end{equation}
With associative identities and distributivity, the invariant classes form a semiring-like algebra. Stabilized idempotent commuting invariants with complements produce a Boolean-like classical regime. Destructive recombination of unresolved alternatives requires a richer signed or phased valuation before stabilization.

\section{Quantum representation programme}
Four explicit principles are used:
\begin{enumerate}
\item reenactable formation: consistent refinement and recombination preserve the invariant;
\item standpoint reciprocity: pure active formations and pure stabilized distinctions can serve as each other's operational seats;
\item continuous reversible standpoint symmetry: equal-weight pure formations are connected by continuous reversible transformations;
\item relational completeness: a composite is determined by the full accessible relational statistics of its components.
\end{enumerate}
Under spectral decomposability, homogeneity, self-duality, finite dimension, and standard composition hypotheses, established representation theorems yield Euclidean Jordan structure. Local tomography and a nonclassical qubit-like sector select the finite-dimensional complex quantum branch.

\claimstatus{Conditional representation. The four principles are not proved minimal or derivable from the primitive formation alone.}

The represented objects are positive Hermitian states, effects, trace probabilities, unitary reversible dynamics, and tensor-product composition. The numerical value of $\hbar$ is a physical calibration, not derived by the formal calculus.

\section{Causal quotient and Lorentzian representation}
Continuation defines a partial order. Incomparability is the primitive analogue of spacelike separation. Under strong causality and smooth-continuum regularity, causal order determines conformal Lorentzian structure; a locally finite occurrence measure fixes scale. Proper time is recovered from calibrated chains.

\claimstatus{Conditional representation. Existence and convergence of a smooth continuum from arbitrary finite formation networks remain open.}

\section{Finite unified loop}
A finite causal diamond supplies one active phase projection and one geometric transport projection. In a minimal abelian model, an edge connection $A_e=\chi_e+i\theta_e$ yields loop holonomy
\begin{equation}
H=A_{OL}+A_{LF}-A_{RF}-A_{OR}.
\end{equation}
The imaginary projection gives interference probabilities, the real projection gives a Lorentz boost rapidity, and adjacent cells compose with cancellation of the shared internal edge. This is a proof of compositional consistency for the toy representation, not a theory of four-dimensional gravity.

\section{Block connection and curvature}
Introduce the block connection
\begin{equation}
\mathcal A=\begin{pmatrix}\Omega&\Upsilon\\\widetilde\Upsilon&\Theta\end{pmatrix},
\end{equation}
with curvature
\begin{equation}
\mathcal F=d\mathcal A+\mathcal A\wedge\mathcal A
=\begin{pmatrix}
R_\Omega+\Upsilon\wedge\widetilde\Upsilon&D\Upsilon\\
D\widetilde\Upsilon&R_\Theta+\widetilde\Upsilon\wedge\Upsilon
\end{pmatrix}.
\end{equation}
If $\Upsilon=\eps U$ and $\widetilde\Upsilon=\eps V$, off-diagonal conversion curvature is $O(\eps)$ and diagonal backreaction is $O(\eps^2)$.

\begin{theorem}[Trace exchange]
The quadratic conversion terms obey
\begin{equation}
\Tr(\delta R_\Omega+\delta R_\Theta)=0.
\end{equation}
\end{theorem}
\begin{proof}
The diagonal terms are commutator-like products $UV-VU$ in opposite block order. Cyclicity of trace cancels their sum.
\end{proof}

\section{Conversion isospectrality}
For one reciprocal map $U$, the two diagonal positive operators are $UU^\dagger$ and $U^\dagger U$.
\begin{theorem}[Nonzero spectral equality]
$UU^\dagger$ and $U^\dagger U$ have identical nonzero spectra with multiplicity.
\end{theorem}
\begin{proof}
A singular-value decomposition $U=W\Sigma V^\dagger$ gives $UU^\dagger=W\Sigma^2W^\dagger$ and $U^\dagger U=V\Sigma^2V^\dagger$.
\end{proof}
The singular spectrum is a parameter-free algebraic invariant of the reciprocal map. It becomes a physical prediction only after both sectors and the map are defined independently of the measured outcomes.

\section{Retired physical branch}
An unrestricted mixed-path factorization of two equal-trace positive matrices is always possible and therefore nonpredictive. A tested single-path comparison did not support the proposed equality at the supplied point estimate. The physical conversion interpretation is therefore not asserted. The exact block algebra, scaling, trace, and singular-spectrum theorems remain valid.

\section{Field equation and conservation}
A Yang-Mills-type candidate action gives
\begin{equation}
D_\mu\mathcal F^{\mu\nu}=g^2\mathcal J^\nu,
\end{equation}
with covariant source conservation from the Bianchi identity. This is a candidate dynamical class, not a unique consequence of the primitive calculus. Higher-curvature, nonlocal, and additional-field actions are not excluded.

\section{Correspondence limits}
When conversion is negligible and geometry fixed, the active representation reduces to unitary quantum evolution. When alternatives are stabilized and the geometric block is represented by a Lorentz connection and coframe, standard four-dimensional assumptions recover an Einstein-type equation with cosmological term. These are correspondence statements rather than derivations of observed constants.

\section{Open theorem and experiment list}
\begin{enumerate}
\item prove minimality or independence of the formation principles;
\item prove finite-to-continuum convergence;
\item derive a preferred action internally;
\item identify an invariant observable for the conversion block before inspecting data;
\item test that observable on held-out instrument and geometric measurements;
\item reject the physical conversion branch if every effect reduces to ordinary open-system dynamics and semiclassical geometry.
\end{enumerate}


\section{Historical residual and correction}
The original categorical route used a fixed-point-free record twist and diagonal self-application to prove that a current self-model omits the act generated by its own reflexive use. This established a residual but did not by itself create new mathematics below the already supplied objects, exponentials, and record space. The formation calculus changes the primitive from object plus report to an occurrence in which active role, stabilized role, and changed standpoint arise together.

\section{Reenactment fixed point}
A seat $[P]$ is stabilized content capable of reenacting pattern $P$. Let $\Delta(x)$ enact the seat obtained by applying $x$ to itself. For every seat-expressible transformer $T$, choose a seat $q_T$ whose enactment on $x$ returns $T(\Delta(x))$. Then
\begin{equation}
\Omega_T=\Delta(q_T)\simeq T(\Omega_T).
\end{equation}
Taking $T$ to be completion gives a conditional final formation: completion adds no new structural-functional pattern. This is closure of representational form, not an assertion that no new empirical content can arrive.

\section{Generated probes and endogenous completion}
Let $\mathcal P(P)$ be the smallest probe family generated by enactment, stabilization, composition, and successor-state change. The complete signature is
\begin{equation}
\Sigma(P)=\bigl(q(P)\bigr)_{q\in\mathcal P(P)}.
\end{equation}
Formation equivalence is equality of complete signatures. A law is an automorphism preserving all generated probes. Completion adjoins the generated probes and their results. If the probe algebra is closed under composition,
\begin{equation}
K(K(P))\simeq K(P),
\end{equation}
and regenerating the law algebra is likewise idempotent. This removes an external declaration of completeness: closure is tested against the operations generated by the formation itself.

\section{Invariant algebra}
Resonance classes support sequential composition $\odot$ and unresolved-alternative composition $\boxplus$. In stabilized idempotent, commuting, complemented regimes the algebra reduces to a Boolean-like classical structure. Before stabilization, destructive recombination requires a signed or phased valuation richer than nonnegative probability.

\section{Quantum representation assumptions}
The route to finite complex quantum theory uses four explicit properties: spectral/re-enactable formation, state-effect reciprocity, continuous reversible symmetry on pure formations, and relational completeness of composites. Under finite-dimensional regularity these properties motivate homogeneous self-dual cones and Euclidean Jordan structure; standard composition restrictions then select the complex branch. This is a representation theorem programme, not a proof that the four properties follow uniquely from the primitive.

\section{Causal quotient and continuum}
Continuation induces a partial order. Under strong causality and smooth continuum regularity, causal order determines conformal Lorentzian geometry; a locally finite formation count fixes scale. Existence and convergence of that continuum from arbitrary finite formation networks remain open.

\section{Block connection and exact algebra}
Write
\begin{equation}
\mathcal A=\begin{pmatrix}\Omega&\Upsilon\\\widetilde\Upsilon&\Theta\end{pmatrix},
\qquad
\mathcal F=d\mathcal A+\mathcal A\wedge\mathcal A.
\end{equation}
Then
\begin{equation}
\mathcal F=\begin{pmatrix}
R_\Omega+\Upsilon\wedge\widetilde\Upsilon&D\Upsilon\\
D\widetilde\Upsilon&R_\Theta+\widetilde\Upsilon\wedge\Upsilon
\end{pmatrix}.
\end{equation}
If $\Upsilon=\varepsilon U$ and $\widetilde\Upsilon=\varepsilon V$, conversion curvature is first order and diagonal backreaction second order. Cyclicity of trace gives
\begin{equation}
\Tr(\delta R_\Omega+\delta R_\Theta)=0.
\end{equation}
For reciprocal conversion, $UU^\dagger$ and $U^\dagger U$ have identical nonzero spectra. This is an exact singular-value theorem and a parameter-free normalized relation inside the block model.

\section{Retired physical interpretation}
The public four-record instrument did not support the proposed single-map spectral equality at its point estimate. Allowing arbitrary mixed paths after observing both marginals is nonpredictive because every equal-trace positive pair admits such a factorization. The manuscript therefore retains the block algebra and isospectral theorem while retiring the unrestricted physical conversion claim. A future physical proposal must specify an invariant observable and prediction before the comparison data are inspected.

\section{Open theorem programme}
The remaining tasks are: derive or reduce the representation assumptions; prove finite-to-continuum convergence; derive a dynamical action rather than postulating a field equation; identify a conversion observable invariant under representation change; and obtain an independent experimental discriminator beyond standard open-system dynamics and semiclassical backreaction.

\begin{thebibliography}{99}
\bibitem{lawvere} F. W. Lawvere, Diagonal arguments and cartesian closed categories, Lecture Notes in Mathematics 92, 1969.
\bibitem{birkedal} L. Birkedal et al., First steps in synthetic guarded domain theory, Logical Methods in Computer Science 8 (2012).
\bibitem{barnum} H. Barnum and A. Wilce, Local tomography and the Jordan structure of quantum theory, Foundations of Physics 44 (2014), 192--212.
\bibitem{hardy} L. Hardy, Quantum theory from five reasonable axioms, arXiv:quant-ph/0101012.
\bibitem{chiribella} G. Chiribella, G. M. D'Ariano, and P. Perinotti, Informational derivation of quantum theory, Physical Review A 84 (2011), 012311.
\bibitem{malament} D. Malament, The class of continuous timelike curves determines the topology of spacetime, Journal of Mathematical Physics 18 (1977), 1399--1404.
\bibitem{loves} D. Lovelock, The Einstein tensor and its generalizations, Journal of Mathematical Physics 12 (1971), 498--501.
\bibitem{wise} D. K. Wise, MacDowell-Mansouri gravity and Cartan geometry, Classical and Quantum Gravity 27 (2010), 155010.
\end{thebibliography}
\end{document}
